\subsection{Protocolo de Instrumentação, Isolamento e Metodologia Analítica}
\label{sec:instrumentacao}

Para assegurar que os tempos e picos de memória medidos reflitam o comportamento algorítmico
real e não artefatos do ambiente de medição, a camada de instrumentação foi desenhada segundo
três princípios metodológicos:

\begin{enumerate}
  \item \textbf{Separação estrita entre passadas de tempo puro e de rastreamento de memória}:
  Rastreadores finos de alocação de memória em nível de interpretador interceptam cada objeto
  criado e destruído durante a execução. Em algoritmos recursivos que alocam e liberam
  estruturas milhões de vezes por segundo, manter esse rastreamento ativo durante as medições
  temporais introduz um severo \emph{efeito observador}, inflando artificialmente o tempo de
  execução e antecipando de forma espúria o limite de 30 minutos. Por isso, todas as medições
  oficiais de tempo executam com rastreamento fino de alocações desativado (utilizando relógio
  monotônico de alta resolução em nanossegundos e tomando a mediana de múltiplas repetições nos
  algoritmos polinomiais para filtrar oscilações do escalonador do sistema operacional). O pico
  relativo de memória por subrotina é coletado em uma passada dedicada e isolada.

  \item \textbf{Isolamento em subprocessos, memória física do sistema e captura limpa de OOM}:
  Cada caso experimental do \textit{Branch-and-Bound} é disparado em um subprocesso
  independente do sistema operacional. Esse isolamento oferece três garantias essenciais:
  (i) permite consultar diretamente o pico máximo de memória física residente (\textit{Maximum
  Resident Set Size}, Max RSS) reportado pelo núcleo do sistema operacional sem impor qualquer
  sobrecarga de CPU à busca; (ii) impede que fragmentação de memória ou ciclos pendentes do
  coletor de lixo de uma instância contaminem a medição da instância seguinte; e (iii) possibilita
  impor um teto explícito de memória virtual nas execuções em CPU. Quando uma versão baseada em
  fila global atinge esse teto, a exceção de memória é capturada de forma controlada, a fila é
  esvaziada para recuperar espaço operacional e o caso é registrado com estado \texttt{OOM},
  preservando o tamanho exato da fronteira e o total de nós expandidos no instante do colapso.
  Todos os resultados são persistidos imediatamente em disco de maneira atômica para permitir
  retomada segura de baterias longas.

  \item \textbf{Modelagem estatística de complexidade e sensibilidade topológica}:
  Na fase analítica (Seção~\ref{sec:experimentos}), as séries temporais e espaciais obtidas por
  subrotina são ajustadas às curvas teóricas previstas por regressão não linear por mínimos
  quadrados, avaliando a aderência pelo coeficiente de determinação $R^2$ e pela raiz do erro
  quadrático médio (RMSE). Para investigar a sensibilidade topológica das aproximações e do
  limite inferior da busca exata, extraímos oito descritores geométricos e estruturais da AGM
  por instância, modelando sua influência por regressão linear múltipla padronizada com
  regularização L2 (regressão \textit{Ridge} com validação cruzada \textit{Leave-One-Out},
  adequada à presença de colinearidade entre descritores em $n = 15$ instâncias), complementada
  por árvores de decisão rasas interpretáveis e florestas aleatórias para capturar regras não
  lineares e limiares críticos.
\end{enumerate}
